\chapter{Conclusion}
\label{sec:conclusion}

We presented the ANU team's submission to PlantCLEF 2026: a partial fine tune of BioCLIP 2.5 ViT-H/14 with only the last four transformer blocks unfrozen, trained on the enlarged PlantCLEF~2024 and iNaturalist (i001) manifest with taxonomy aware auxiliary heads, and evaluated under tiled softmax mean inference with class prior logit adjustment and a $224{+}336$ resolution ensemble. It reached a public Macro F1 of \textbf{0.41826} and a private Macro F1 of \textbf{0.40283}.

We highlight three findings. Unfreeze depth has a substantial effect on BioCLIP 2.5: a moderate depth of four to eight blocks is the operating region, whereas a full fine tune degrades the Tree-of-Life prior and reduces the model to $0.301$ public. Single plant validation accuracy and quadrat Macro F1 are only weakly related ($r \approx 0.4$, and they diverge within the unfreeze ablation), which indicates that the limiting factor is the distribution shift between single plant and multi species images and that model capacity is already sufficient. Post hoc reweighting of a strong visual posterior adds little: the triple backbone late fusion classifier ($\Delta{=}-0.036$), the genus/family co-occurrence rerank ($\Delta{=}-0.006$), and the seasonal phenology prior ($\Delta{=}-0.005$) all failed to improve the anchor, because fine tuned visual classifiers on biological data absorb such signals implicitly through training distribution co-occurrence. The $010\rightarrow$i002 gain is confounded across data, schedule, and architecture, so we make no causal claim for any single change; \S\ref{sec:discussion} details the limitations and future work.
